\subsection{Planitud y profundidad.}

\begin{proposition}[6.3.1]
\EGAReaderTarget{ega4.c04.prop.6-3-1}{IV.2.s6.3.1-fr}
Sean $A$, $B$ dos anillos locales noetherianos, $k$ el cuerpo residual de
$A$, $\varphi:A\to B$ un homomorfismo local, $M$ un $A$-módulo de tipo finito
y $N$ un $B$-módulo de tipo finito. Si $N$ es un $A$-módulo plano $\neq0$, se
tiene
\begin{equation}
\operatorname{prof}_B(M\otimes_AN)=\operatorname{prof}_A(M)
+\operatorname{prof}_{B\otimes_Ak}(N\otimes_Ak).
\tag{6.3.1.1}\label{IV.2.s6.3.1.1-fr}
\end{equation}
\end{proposition}

Es posible limitarse al caso en que $M\neq0$, pues de lo contrario los dos
miembros de (6.3.1.1) son iguales a $+\infty$. Procederemos por inducción sobre
el entero $n$ igual al segundo miembro de (6.3.1.1) (que es finito en virtud
de las hipótesis, de (\textbf 0, 16.4.6.2) y del lema de Nakayama aplicado a $N$).
Supongamos primero $n=0$; entonces, si $\mathfrak m$ y $\mathfrak n$ designan
los ideales maximales de $A$ y $B$, respectivamente, $\mathfrak m$ es un ideal
primo asociado a $M$ y $\mathfrak n/\mathfrak mB$ un ideal primo de
$B/\mathfrak mB=B\otimes_Ak$ asociado a $N\otimes_Ak$ (\textbf 0, 16.4.6, (i)). Se
concluye entonces de (3.3.1) que $\mathfrak n$ es un ideal primo de $B$
asociado a $M\otimes_AN$ y, por tanto, (\textbf 0, 16.4.6, (i)), que
$\operatorname{prof}_B(M\otimes_AN)=0$. Supongamos, pues, $n>0$, y distingamos
dos casos:

$a)$ Supongamos $\operatorname{prof}_A(M)>0$. Sea $x\in\mathfrak m$ un
elemento $M$-regular, y pongamos
\[
A'=A/xA,\qquad B'=B/xB,\qquad M'=M/xM,\qquad N'=N/xN;
\]
se tiene
\[
B'\otimes_{A'}k=B\otimes_Ak,\qquad N'\otimes_{A'}k=N\otimes_Ak
\]
puesto que $x\in\mathfrak m$, y
\[
M'\otimes_{A'}N'=(M\otimes_AN)/x(M\otimes_AN);
\]

\oldpage[IV-2]{139}

además, como $N$ es un $A$-módulo plano, la hipótesis de que $x$ es
$M$-regular implica que $x$ es también $(M\otimes_AN)$-regular
$(\mathbf{0}_{\mathrm I}, 6.1.1)$. Se tiene, por tanto, (\textbf 0, 16.4.6, (ii) y 16.4.8)
\[
\operatorname{prof}_{A'}(M')=\operatorname{prof}_A(M)\setmin 1,
\qquad
\operatorname{prof}_{B'}(M'\otimes_{A'}N')
=\operatorname{prof}_B(M\otimes_AN)\setmin 1
\]
y
\[
\operatorname{prof}_{B'\otimes_{A'}k}(N'\otimes_{A'}k)
=\operatorname{prof}_{B\otimes_Ak}(N\otimes_Ak).
\]

La igualdad (6.3.1.1) es, pues, consecuencia de la misma relación para $A'$,
$B'$, $M'$ y $N'$; pero, como $N$ es un $A$-módulo plano,
$N'=N\otimes_AA'$ es un $A'$-módulo plano $(\mathbf{0}_{\mathrm I}, 6.2.1)$; se puede,
por tanto, aplicar la hipótesis de inducción, lo que demuestra (6.3.1.1) en
este caso.

$b)$ Supongamos $\operatorname{prof}_{B\otimes_Ak}(N\otimes_Ak)>0$.
Consideremos un elemento $y\in\mathfrak m$ que sea
$(N\otimes_Ak)$-regular; resulta de $(\mathbf{0}_{\mathrm{III}}, 10.2.4)$ que $y$ es
entonces $N$-regular y que
\[
N'=N/yN
\]
es un $A$-módulo \emph{plano}, puesto que se supone que $N$ es un $A$-módulo plano.
Aplicando $(\mathbf{0}_{\mathrm I}, 6.1.2)$ a la sucesión exacta de $A$-módulos
\[
0\to N\xrightarrow{y}N\to N'\to0
\]
se concluyen los isomorfismos
\[
(M\otimes_AN)/y(M\otimes_AN)\xrightarrow{\sim}M\otimes_AN'
\quad\text{y}\quad
N'\otimes_Ak\xrightarrow{\sim}(N\otimes_Ak)/y(N\otimes_Ak)
\]
y, además, que $y$ es $(M\otimes_AN)$-regular. Se tiene, por tanto,
\[
\operatorname{prof}_{B\otimes_Ak}(N'\otimes_Ak)
=\operatorname{prof}_{B\otimes_Ak}(N\otimes_Ak)\setmin 1
\]
y
\[
\operatorname{prof}_B(M\otimes_AN')
=\operatorname{prof}_B(M\otimes_AN)\setmin 1;
\]
la hipótesis de inducción muestra que la relación (6.3.1.1) es válida para
$A$, $B$, $M$ y $N'$, y de lo que precede se deduce (6.3.1.1) para $A$, $B$,
$M$ y $N$.
